\begin{appendices}
\section{Rastreabilidade e auditoria da base experimental}\label{ap:enquadramento}
\begin{fillblock}
\noindent\textbf{[A PREENCHER]} Incluir tabela com ID e nomes da Tabela 1 de \textcite{Wolenski2020}, médias e DP das Tabelas 3--6, resistências características e classe publicada da Tabela 7, classe recalculada e classe obtida pela conversão simplificada. Registrar a tabela/página de origem, divergências e qualquer tratamento efetuado. Separar os valores transcritos dos calculados e das hipóteses de material. O módulo $E_{M0}$ não deve aparecer como dado medido da fonte.
\end{fillblock}

\section{Memorial de verificação e reprodutibilidade}\label{ap:relatorio}
\begin{fillblock}
\noindent\textbf{[A PREENCHER]} Apresentar as entradas e a conferência manual de um caso sob geometria fixa, seguida dos projetos R e C e da reavaliação C$\rightarrow$R. Documentar coeficientes, unidades, hipóteses complementares, sementes, critérios de convergência e tolerâncias de viabilidade. Disponibilizar os dados derivados e os scripts usados nas tabelas e figuras.
\end{fillblock}
\end{appendices}
